% Gabarit pandoc des corrigés publiés (brut et guidé). Contenu sans marque ni nom de site.
% Compilation : lualatex (latexmk -lualatex).
% Le Markdown reste la source unique : ne pas éditer ce fichier .tex à la main.
\documentclass[11pt,a4paper]{article}

% --- Moteur, polices (Latin Modern) et langue ---------------------------------
\usepackage{iftex}
\ifLuaTeX\else\errmessage{Ce document se compile avec lualatex}\fi
\usepackage{amsmath,amssymb,mathtools}
\usepackage{fontspec}
% mathrm=sym : \mathrm prend ses glyphes dans la police mathématique (ℓ, μ, Ω), comme KaTeX.
\usepackage[math-style=TeX,bold-style=TeX,mathrm=sym]{unicode-math}
% Repli pour les glyphes absents de Latin Modern (indices et exposants Unicode,
% flèches, filets des schémas en texte, symboles) : DejaVu, puis Latin Modern Math
% (ℓ, lettres mathématiques Unicode des légendes).
\directlua{
  luaotfload.add_fallback("repliTexte", {"DejaVuSerif:mode=harf;", "DejaVuSans:mode=harf;", "file:latinmodern-math.otf:mode=harf;"})
  luaotfload.add_fallback("repliMono", {"DejaVuSansMono:mode=harf;", "DejaVuSans:mode=harf;", "file:latinmodern-math.otf:mode=harf;"})
}
\setmainfont{Latin Modern Roman}[RawFeature={fallback=repliTexte}]
\setsansfont{Latin Modern Sans}[RawFeature={fallback=repliTexte}]
\setmonofont{Latin Modern Mono}[RawFeature={fallback=repliMono}]
\setmathfont{Latin Modern Math}
\usepackage[french]{babel}
\frenchsetup{SmallCapsFigTabCaptions=false}

% --- Mise en page --------------------------------------------------------------
\usepackage[a4paper,margin=2.2cm,headheight=14pt,footskip=1.2cm]{geometry}
\usepackage{microtype}
\setlength{\emergencystretch}{4em}
\tolerance=1500
\setlength{\parindent}{0pt}
\setlength{\parskip}{0.55\baselineskip plus 2pt minus 1pt}
\usepackage{xcolor}
% Gris nommés : babel-french rend « ! » actif dans le document.
\definecolor{gris2}{gray}{0.98}
\definecolor{gris3}{gray}{0.97}
\definecolor{gris5}{gray}{0.95}
\definecolor{gris35}{gray}{0.65}
\definecolor{gris40}{gray}{0.60}
\definecolor{gris45}{gray}{0.55}
\definecolor{gris50}{gray}{0.50}
\definecolor{gris55}{gray}{0.45}
\definecolor{gris60}{gray}{0.40}
\definecolor{gris65}{gray}{0.35}
\definecolor{gris70}{gray}{0.30}
\definecolor{gris80}{gray}{0.20}
\definecolor{gris85}{gray}{0.15}
\usepackage{graphicx}
\usepackage{float}
\usepackage{caption}
\captionsetup{font=small,labelformat=empty,justification=centering}
\usepackage{booktabs,longtable,array,multirow,calc}
\usepackage{siunitx}
\sisetup{locale=FR}
\usepackage[version=4]{mhchem}
\usepackage{listings}
\usepackage{fancyhdr}
\usepackage{lastpage}
\usepackage[most]{tcolorbox}
\usepackage{enumitem}
\setlist{nosep,leftmargin=*,topsep=2pt}
\usepackage[hidelinks,unicode]{hyperref}
\usepackage{bookmark}
\hypersetup{pdftitle={G2E BCPST — Composition française 2026},pdfauthor={Rédigé par IA},pdfcreator={LaTeX},pdfsubject={Corrigé brut}}
\usepackage{xurl}

% --- Pandoc : listes, tableaux, figures ----------------------------------------
\providecommand{\passthrough}[1]{#1}
\providecommand{\tightlist}{\setlength{\itemsep}{0pt}\setlength{\parskip}{0pt}}

\makeatletter
\def\maxwidth{\ifdim\Gin@nat@width>\linewidth\linewidth\else\Gin@nat@width\fi}
\def\maxheight{\ifdim\Gin@nat@height>0.55\textheight 0.55\textheight\else\Gin@nat@height\fi}
\makeatother
\setkeys{Gin}{width=\maxwidth,height=\maxheight,keepaspectratio}
\setcounter{secnumdepth}{-\maxdimen}
\providecommand{\pandocbounded}[1]{#1}

% --- Code (listings) -----------------------------------------------------------
\lstset{
  basicstyle=\ttfamily\small,
  columns=fullflexible,
  keepspaces=true,
  breaklines=true,
  breakatwhitespace=false,
  postbreak=\mbox{\textcolor{gris50}{$\hookrightarrow$}\space},
  frame=single,
  rulecolor=\color{gris35},
  framesep=4pt,
  xleftmargin=5pt,xrightmargin=5pt,
  backgroundcolor=\color{gris3},
  keywordstyle=\bfseries,
  commentstyle=\color{gris65},
  stringstyle=\color{gris80},
  showstringspaces=false,
  tabsize=4,
  extendedchars=true,
  aboveskip=6pt,belowskip=6pt,
}

% --- Macros des corrigés (KaTeX ↔ LaTeX) ---------------------------------------
% Les corrigés sont écrits pour KaTeX : on définit ici ce que KaTeX connaît et
% que LaTeX n'a pas, sans jamais retoucher le Markdown.
\providecommand{\R}{\mathbb{R}}
\providecommand{\N}{\mathbb{N}}
\providecommand{\Z}{\mathbb{Z}}
\providecommand{\Q}{\mathbb{Q}}
\providecommand{\C}{\mathbb{C}}
\providecommand{\dd}{\mathrm{d}}
\providecommand{\llbracket}{\lBrack}
\providecommand{\rrbracket}{\rBrack}
\providecommand{\lBrace}{\lbrace}
\providecommand{\rBrace}{\rbrace}
\providecommand{\Box}{\square}
\providecommand{\Diamond}{\lozenge}
\providecommand{\varnothing}{\emptyset}
\providecommand{\argmax}{\operatorname*{arg\,max}}
\providecommand{\argmin}{\operatorname*{arg\,min}}
\renewcommand{\leq}{\leqslant}
\renewcommand{\geq}{\geqslant}

% Équation hors texte : réduite à la largeur de la ligne si elle déborde.
\newsavebox{\boiteequation}
\newcommand{\equationajustee}[1]{%
  \sbox{\boiteequation}{$\displaystyle #1$}%
  \ifdim\wd\boiteequation>\linewidth
    \[\resizebox{\linewidth}{!}{\usebox{\boiteequation}}\]%
  \else
    \[#1\]%
  \fi}

% --- Encadrés pédagogiques (sobres, lisibles en noir et blanc) -----------------
\tcbset{
  encadre/.style={
    enhanced,breakable,sharp corners,
    colback=white,colframe=gris55,
    boxrule=0.4pt,leftrule=2.2pt,
    left=6pt,right=6pt,top=3pt,bottom=3pt,
    before skip=6pt,after skip=6pt,
    fonttitle=\bfseries\small,coltitle=black,
    attach title to upper,after title={\quad},
    parbox=false,
  }
}
\newtcolorbox{indication}[1]{encadre,colframe=gris45,colback=gris2,title={Indication~#1}}
\newtcolorbox{methode}{encadre,colframe=gris70,title={Méthode}}
\newtcolorbox{piege}{encadre,colframe=gris85,boxrule=0.8pt,leftrule=2.2pt,title={Piège}}
\newtcolorbox{jury}{encadre,colframe=gris60,colback=gris5,title={Jury}}
\newtcolorbox{encadrenote}{encadre,colframe=gris40}

% --- Titres de section ---------------------------------------------------------
\usepackage{titlesec}
\titleformat{\section}{\Large\bfseries}{}{0pt}{}
\titleformat{\subsection}{\large\bfseries}{}{0pt}{}
\titleformat{\subsubsection}{\normalsize\bfseries}{}{0pt}{}
\titlespacing*{\section}{0pt}{14pt plus 4pt minus 2pt}{6pt}
\titlespacing*{\subsection}{0pt}{10pt plus 3pt minus 2pt}{4pt}
\titlespacing*{\subsubsection}{0pt}{8pt plus 2pt minus 2pt}{2pt}

% --- En-tête et pied de page ---------------------------------------------------
\pagestyle{fancy}
\fancyhf{}
\fancyhead[L]{\small\textsc{Corrigé brut} · G2E BCPST}
\fancyhead[R]{\small Composition française 2026}
\fancyfoot[L]{\footnotesize Corrigé rédigé par IA, non officiel}
\fancyfoot[R]{\footnotesize \thepage\,/\,\pageref*{LastPage}}
\renewcommand{\headrulewidth}{0.4pt}
\renewcommand{\footrulewidth}{0pt}
\fancypagestyle{plain}{\fancyhf{}\fancyfoot[R]{\footnotesize \thepage\,/\,\pageref*{LastPage}}\fancyfoot[L]{\footnotesize Corrigé rédigé par IA, non officiel}\renewcommand{\headrulewidth}{0pt}}

\begin{document}
\thispagestyle{plain}
\begin{center}
  {\small\scshape Corrigé brut}\par\smallskip
  {\LARGE\bfseries G2E BCPST — Composition française 2026\par}\medskip
  {\large G2E · BCPST · Composition française · 2026\par}
\end{center}
\vspace{2pt}\hrule\vspace{8pt}

\emph{Corrigé rédigé par IA, non officiel. Signalez toute erreur.}

Épreuve de 3 h 30, sans document. Sujet : confronter un propos de Jean Giono (fin des \emph{Vraies richesses}) aux trois œuvres au programme, Jules Verne, \emph{Vingt mille lieues sous les mers} ; Georges Canguilhem, \emph{La Connaissance de la vie} ; Marlen Haushofer, \emph{Le Mur invisible}. L'épreuve comporte une seule question, la dissertation ; aucune limite de mots n'est imposée. Les citations des œuvres sont données de mémoire ; celles dont la lettre n'est pas assurée sont rapportées au style indirect.

\hypertarget{dissertation}{%
\section{Dissertation}\label{dissertation}}

\textbf{Problématique :} la richesse du monde vivant est-elle un bien qui nous était destiné, qu'il suffirait de reprendre et où la mort viendrait se fondre, ou une relation à instituer, par le travail et la médiation, avec un monde qui ne nous attendait pas ?

\textbf{Plan :} I. Un diagnostic partagé : les œuvres écrivent contre une dépossession. II. Rien n'a été « fait pour » l'homme : un monde indifférent, que l'on n'atteint que par médiation. III. Reste l'impératif : une richesse à instituer, que la mort n'accomplit pas.

\hypertarget{introduction}{%
\subsection{Introduction}\label{introduction}}

Est riche, croit-on, celui qui possède. Jean Giono retourne le mot contre l'économie qui l'a confisqué. Au terme des \emph{Vraies richesses}, il s'adresse à un « tu » qui est chacun de nous : « Ce dont on te prive, c'est de vents, de pluies, de neiges, de soleils, de montagnes, de fleuves, et de forêts : les vraies richesses de l'homme ! Tout a été fait pour toi ; au fond de tes plus obscures veines, tu as été fait pour tout. Quand la mort arrivera, ne t'inquiète pas, c'est la continuation logique. Tâche seulement d'être le plus riche possible. A ce moment-là, ce que tu es, deviens. »

Le propos enchaîne quatre affirmations. Il dénonce une spoliation dont l'auteur reste dans l'ombre d'un « on » : la ville, l'argent, la machine, tout ce qui s'interpose. Il dresse l'inventaire des biens dérobés, qui ne sont ni des objets ni des œuvres mais des éléments, tous au pluriel, inutiles et inappropriables : on ne fait rien d'un vent, on ne le thésaurise pas. Il fonde ensuite ce grief sur un chiasme, « tout a été fait pour toi {[}\ldots{]} tu as été fait pour tout », qui affirme une convenance réciproque de l'homme et du monde jusque dans la part organique et inconsciente de nous-mêmes, ces « obscures veines ». Il en tire enfin une sagesse de la mort : puisque nous sommes de la même étoffe que le monde, mourir n'est qu'une « continuation logique », et l'unique tâche est d'être « le plus riche possible », car « à ce moment-là », c'est-à-dire à l'instant de la mort, chacun devient ce qu'il est. Reprenant l'injonction de Pindare, Giono fait de la mort un accomplissement : elle rend au tout, intégralement, ce que la vie a reçu de lui.

Le texte recèle pourtant une tension. Si nous avons été faits pour tout, comment pouvons-nous en être privés ? Et si la richesse nous est due, pourquoi faut-il « tâcher » de l'acquérir ? Confronter ce propos à trois œuvres qui ont fait du monde vivant leur objet, c'est éprouver ces affirmations une à une. \emph{Vingt mille lieues sous les mers} (1869) est le roman d'un homme qui a quitté la civilisation pour vivre de la mer. \emph{La Connaissance de la vie} (1952) de Georges Canguilhem se demande à quelles conditions on connaît le vivant sans le trahir. \emph{Le Mur invisible} (1963) de Marlen Haushofer rend de force une femme aux neiges et aux forêts, derrière une paroi qui la sépare des hommes.

D'où la question : la richesse du monde vivant est-elle un bien qui nous était destiné, qu'il suffirait de reprendre et dans lequel la mort elle-même viendrait se fondre, ou une relation qu'il faut instituer, au prix d'un travail et d'une médiation, avec un monde qui ne nous attendait pas ?

Les trois œuvres partagent d'abord le diagnostic de Giono : elles naissent d'une séparation et opposent aux fausses richesses celles du vivant. Aucune, cependant, ne souscrit au chiasme : le monde n'a pas été fait pour l'homme, il lui est indifférent et ne se donne qu'à travers des instruments. Reste l'impératif, qui survit à l'épreuve à condition d'y entendre l'institution patiente d'un lien, que la mort interrompt au lieu de l'accomplir.

\hypertarget{i.-un-diagnostic-partaguxe9-les-ux153uvres-uxe9crivent-contre-une-duxe9possession}{%
\subsection{I. Un diagnostic partagé : les œuvres écrivent contre une dépossession}\label{i.-un-diagnostic-partaguxe9-les-ux153uvres-uxe9crivent-contre-une-duxe9possession}}

\textbf{A. Trois œuvres nées d'une privation}

Aucune des trois œuvres ne s'écrit depuis une plénitude ; toutes s'écrivent depuis une séparation, et c'est d'abord en cela qu'elles donnent raison à Giono.

Chez Verne, la séparation est un programme. Nemo ne s'est pas vu privé de la mer, il s'est retranché de la terre : « Je ne suis pas ce que vous appelez un homme civilisé ! J'ai rompu avec la société tout entière », déclare-t-il à Aronnax. Ce que la société lui refusait, c'est la liberté, et cette liberté a un lieu : « La mer n'appartient pas aux despotes. » À sa surface, les puissants peuvent encore se battre et s'entre-dévorer, mais à trente pieds au-dessous, « leur pouvoir cesse, leur influence s'éteint, leur puissance disparaît ». Le « on » anonyme de Giono reçoit ici un nom politique : les empires qui oppriment les peuples.

Chez Haushofer, la privation est subie et totale. Un matin, la narratrice se heurte à une paroi transparente et froide qui ferme la vallée ; au-delà, un vieil homme est pétrifié au bord d'une fontaine. L'événement n'est jamais expliqué, et il n'a pas besoin de l'être pour agir : il soustrait la radio, les voisins, les nouvelles, l'usage de la voiture. Or la narratrice découvre, non sans stupeur, que sa vie d'avant, celle d'une veuve occupée de conventions et de petites obligations, ne lui manque presque pas. Elle avait été privée de quelque chose sans le savoir : c'est le mur qui le lui révèle.

Chez Canguilhem, la privation est intellectuelle, et elle est la difficulté même que le livre affronte. Celui qui veut connaître la vie est un vivant, et il doit pourtant faire de ce qu'il est un objet. Or, écrit-il dans l'introduction, « la vie est formation de formes, la connaissance est analyse des matières informées ». La science, qui décompose et mesure, manque par construction ce qui fait la vie : l'acte même de former. Le biologiste est séparé du vivant par le savoir qui devait l'en rapprocher, et tout le livre cherche les conditions d'un retour : la pensée du vivant, dit Canguilhem, doit tenir du vivant l'idée du vivant.

\textbf{B. Contre l'or, l'inventaire des vraies richesses}

Le mot décisif du propos est « richesses », pris au sens économique pour être retourné. Deux des œuvres pratiquent ce retournement à la lettre.

Verne met en scène l'or lui-même. Dans la baie de Vigo gisent les galions espagnols coulés en 1702 ; les hommes du \emph{Nautilus}, en scaphandre, y ramassent l'or comme on ramasse du bois mort. L'or est au fond, à la place du rebut : il est ce que la mer a englouti, et la mer continue. Sa seule valeur est celle que Nemo lui rend en le destinant aux peuples qui luttent. La vraie richesse est ailleurs, et Nemo la nomme : « La mer est tout ! {[}\ldots{]} Son souffle est pur et sain. C'est l'immense désert où l'homme n'est jamais seul, car il sent frémir la vie à ses côtés. » Même énumération élémentaire que chez Giono, même équivalence entre abondance naturelle et plénitude morale. L'épisode des pêcheries de Ceylan oppose d'ailleurs les deux richesses : Nemo laisse croître dans sa coquille une perle prodigieuse qu'il ne vendra jamais, et il offre un sac de perles au pêcheur indien qu'il vient d'arracher au requin, parce que cet homme est, dit-il en substance, du pays des opprimés, qui est aussi le sien. La perle est richesse quand elle est vie qui croît ou don ; elle est misère quand elle devient marchandise.

Chez Haushofer, le renversement s'opère par la seule force des choses. Dans la vallée close, les valeurs se recomposent : ce qui compte, c'est le foin rentré avant la pluie, le lait de Bella, les pommes de terre de semence, une boîte d'allumettes, une cartouche. L'argent n'a plus personne à qui s'échanger. Et le récit accumule des instants de bonheur d'une intensité qui surprend la narratrice elle-même : l'été à l'alpage, la lumière sur le pré, la présence chaude du chien Lynx. Ce sont terme pour terme les richesses de Giono, et Haushofer leur accorde ce privilège : elles seules résistent quand tout le reste est supprimé.

Canguilhem donne à la thèse son meilleur argument. Dans « Aspects du vitalisme », il refuse de traiter le vitalisme comme une erreur préscientifique ; il y voit l'expression d'une « exigence permanente de la vie dans le vivant ». Le vivant ne subit pas son milieu : il le trie, préfère, repousse ; il pose des valeurs par le seul fait de vivre. C'est exactement ce que Giono appelle richesse, non une quantité de biens, mais l'intensité d'un rapport évaluatif au monde. Là où l'écrivain proclame, l'épistémologue établit : la valeur n'est pas plaquée par l'homme sur une nature neutre, elle apparaît avec le vivant.

\textbf{C. Une richesse d'abord sensible}

Les richesses énumérées par Giono sont toutes des objets de sensation. Les trois œuvres accordent au sensible une priorité que leur genre ne laissait pas attendre.

Le \emph{Nautilus} est une machine paradoxale : chef-d'œuvre d'ingénieur, il sert surtout à rendre la mer visible. Les panneaux du salon s'ouvrent, l'électricité illumine les eaux, et l'inventaire savant des espèces bascule dans le ravissement : la forêt sous-marine de l'île Crespo, les prairies de corail, les phosphorescences. Nemo ne dit pas à Aronnax « étudiez la mer », il lui dit « vivez au sein des mers ». Le savoir n'est pas nié, mais il vient après l'émerveillement.

Chez Haushofer, le sensible gagne sur le langage. La narratrice parle de moins en moins, et sa vie se recompose autour de perceptions : le froid, l'humidité, la fatigue des reins après la fauche, le souffle de la vache dans l'étable. Elle écrit son récit parce qu'elle sent les mots se retirer, et l'on peut lire tout le livre comme la trace d'une existence qui redevient sensation.

Canguilhem fournit la justification de cette priorité. Dans « Le vivant et son milieu », il rappelle que l'univers de la physique, homogène et sans centre, est une construction seconde. Ce qui est premier pour tout vivant, homme compris, c'est un milieu propre, découpé par ses organes et polarisé par ses besoins : « Vivre, c'est rayonner, c'est organiser le milieu à partir d'un centre de référence. » L'expérience du vent et du soleil n'est donc pas une version naïve de la météorologie ; elle est le sol dont la météorologie s'est détachée. Giono a raison sur ce point, et pour de meilleures raisons que les siennes.

Sur le diagnostic, l'accord est large : séparation, renversement des richesses, primauté du sensible. Mais il porte sur la première phrase du propos, la plus facile à admettre. Reste le chiasme, qui fonde le grief sur une convenance réciproque de l'homme et du monde. C'est là que les trois œuvres cessent de suivre Giono.

\hypertarget{ii.-rien-na-uxe9tuxe9-fait-pour-lhomme-un-monde-indiffuxe9rent-que-lon-natteint-que-par-muxe9diation}{%
\subsection{II. Rien n'a été « fait pour » l'homme : un monde indifférent, que l'on n'atteint que par médiation}\label{ii.-rien-na-uxe9tuxe9-fait-pour-lhomme-un-monde-indiffuxe9rent-que-lon-natteint-que-par-muxe9diation}}

\textbf{A. Le monde vivant ne veut rien de nous}

Si tout a été fait pour moi et moi pour tout, le monde est habitable de plein droit, et il suffit d'y rentrer. Les œuvres opposent à cette confiance des démentis qui n'ont rien d'accidentel.

Le roman de l'abondance marine est aussi un roman d'agonies. Pris sous la banquise au retour du pôle Sud, le \emph{Nautilus} voit son air se raréfier, et Verne décrit l'asphyxie dans le détail : les tempes qui battent, les hommes qui se relaient pour entamer la glace. La glace se referme selon des lois qui ignorent Aronnax. Le poulpe géant arrache un homme de l'équipage et l'emporte sans que personne y puisse rien. La violence n'est d'ailleurs pas seulement subie : Nemo lance l'éperon du \emph{Nautilus} contre les cachalots, qu'il juge malfaisants, et le récit d'Aronnax devient celui d'un carnage. La mer n'est pas une communion ; on y mange et on y est mangé, et le contemplateur lui-même tue.

\emph{Le Mur invisible} pousse plus loin le démenti, parce que la narratrice ne peut pas remonter à bord. La vallée n'est pas un paradis : il faut faucher, rentrer le bois, compter les réserves, aider Bella à vêler sans vétérinaire, traverser des hivers. Les épreuves les plus dures sont des indifférences plutôt que des attaques : la neige qui tombe trop tôt, la maladie qui vient sans raison, les chatons perdus. La forêt, comprend la narratrice, ne lui veut ni bien ni mal. Le monde n'est pas hostile, ce qui serait encore une relation ; il est sans destinataire. Et ce n'est pas la nature qui a dressé le mur : la narratrice y voit, vraisemblablement, l'effet d'une arme. La catastrophe qui la rend aux neiges est le produit de cette puissance technique dont Giono déplorait qu'elle nous sépare du monde.

Canguilhem donne à ce démenti son concept. La santé n'est pas une harmonie préétablie avec le milieu ; elle est une marge, une latitude, la capacité de supporter les variations d'un milieu qui ne tient pas ses promesses. Il n'y a pas d'ajustement donné une fois pour toutes, mais un jeu qu'il faut sans cesse reconquérir. Et la vie peut échouer : dans « La monstruosité et le monstrueux », le monstre n'est pas une aberration venue du dehors, il est produit par la vie elle-même ; il est « le vivant de valeur négative ». Une nature faite pour nous ne produirait pas de monstres.

\textbf{B. La finalité est dans le vivant, non dans le monde}

Giono ne dit pas seulement que le monde est bon ; il dit qu'il a été \emph{fait} pour nous et nous pour lui. Même si l'on y entend une participation cosmique plutôt qu'un plan providentiel, la proposition affirme une finalité inscrite dans le monde, et \emph{La Connaissance de la vie} permet de l'examiner comme telle.

« Le vivant et son milieu » retrace la généalogie de la notion. Venue de la mécanique newtonienne, où le fluide ambiant transmet l'action, elle passe en biologie avec Lamarck et Comte et devient le décor auquel les êtres doivent s'adapter. On aboutit à une alternative : ou bien le milieu fabrique le vivant, ou bien il a été disposé pour lui, et c'est le finalisme dont Giono hérite. Canguilhem sort de l'alternative. Il n'existe pas de milieu absolu, le même pour tous ; il y a autant de milieux propres que de vivants, chacun découpé par une activité. Le monde de la tique, que Canguilhem reprend à Uexküll, n'est pas un fragment appauvri du nôtre : c'est un autre monde, organisé par trois signaux et par eux seuls. La question « le monde a-t-il été fait pour l'homme ? » perd alors son objet : il n'y a pas de monde disponible en soi dont on pourrait se demander à qui il était destiné.

Cette critique ne restaure pas pour autant un mécanisme sans fins. La finalité est rapatriée dans l'organisme. C'est le sens de « Machine et organisme », qui renverse le rapport traditionnel : au lieu de comprendre l'organisme sur le modèle de la machine, dont le plan et la fin sont extérieurs, dans la tête de l'ingénieur, il faut comprendre la machine comme un prolongement de l'organisme, qui se répare, se reproduit et tolère des écarts qu'aucune machine ne tolère. Giono a donc raison de refuser un monde inerte et sans valeur, mais il place la finalité au mauvais endroit : elle n'est pas dans un plan qui aurait prévu notre venue, elle est en nous comme en toute bête et toute plante, et elle ne nous garantit rien.

La formule de Nemo, « l'homme n'est jamais seul, car il sent frémir la vie à ses côtés », apparaît alors pour ce qu'elle est : la description d'un milieu propre. C'est le \emph{Nautilus} qui fait de l'océan un monde habitable pour Aronnax. Et la narratrice de Haushofer apprend que la vallée est le milieu du chevreuil, de la corneille et du chat, où elle n'entre qu'en se transformant.

\textbf{C. Il n'y a de contact qu'appareillé : la vitre, la faux, le concept}

De là un dernier démenti, le plus concret. Giono rêve d'immédiateté : se tenir dans le vent, reprendre ce qu'on nous a pris. Les œuvres montrent qu'on n'atteint jamais le vivant les mains nues.

Le roman de l'immersion totale est un roman de l'enfermement. Aronnax passe l'essentiel du voyage dans une coque d'acier, respire un air fabriqué et regarde la mer à travers des vitres épaisses. Quand il sort, c'est en scaphandre, avec un appareil Rouquayrol, une lampe Ruhmkorff et un fusil à air comprimé. La devise du navire, \emph{mobilis in mobili}, dit cette condition : on n'est mobile dans l'élément mobile qu'en emportant un petit monde stable. Le \emph{Nautilus} est un milieu propre construit, sans lequel la richesse de la mer serait inaccessible à un homme.

Chez Haushofer, l'appareillage devient comptabilité. Le récit tient de l'inventaire : combien de cartouches, d'allumettes, de pommes de terre de semence. La narratrice ne vit pas de la nature seule ; elle vit d'un stock hérité de la civilisation perdue, qu'elle prolonge par le travail de la faux, de la houe et du fusil. Aux pluriels inépuisables de Giono répond un monde fini, où rien de ce qui est consommé ne sera remplacé. La richesse n'est pas une abondance offerte ; elle est une décroissance gérée.

Canguilhem étend la médiation à la connaissance. « L'expérimentation en biologie animale » montre que la spécificité des espèces interdit l'extrapolation facile, que l'individualité résiste à la standardisation et qu'isoler un organe, c'est déjà le placer dans un milieu qui n'est pas le sien. « La théorie cellulaire » établit que les concepts par lesquels nous saisissons le vivant charrient des valeurs sociales : la cellule comme individu, l'organisme comme république de cellules. Il n'y a pas plus de regard nu sur la vie que d'immersion nue dans la mer ; entre le vivant et nous, il y a toujours une vitre, un outil ou un concept. La question n'est pas de supprimer la médiation, mais de savoir laquelle reste fidèle à son objet.

Le chiasme ne résiste donc pas : rien n'a été fait pour nous, le milieu est infidèle et nous n'y entrons qu'équipés. Faut-il abandonner tout le propos ? Ce serait oublier qu'il se termine sur un impératif, « tâche seulement d'être le plus riche possible », dont le verbe avoue un effort. C'est cette éthique que les œuvres reprennent, à condition d'en réviser le terme : la mort.

\hypertarget{iii.-reste-limpuxe9ratif-une-richesse-uxe0-instituer-que-la-mort-naccomplit-pas}{%
\subsection{III. Reste l'impératif : une richesse à instituer, que la mort n'accomplit pas}\label{iii.-reste-limpuxe9ratif-une-richesse-uxe0-instituer-que-la-mort-naccomplit-pas}}

\textbf{A. « Tâche » : la richesse comme travail et comme invention de normes}

« Tâcher », c'est s'efforcer sans être assuré du résultat. Le mot corrige de l'intérieur la plainte initiale : une richesse qu'il faut tâcher d'atteindre n'est pas un bien dont on aurait été dépouillé, c'est une conquête. Les trois œuvres l'entendent ainsi.

La narratrice du \emph{Mur invisible} ne retrouve pas une nature perdue : elle devient quelqu'un d'autre, et le récit enregistre cette mue. Les gestes de la vie sociale tombent l'un après l'autre ; elle constate qu'elle n'est plus une femme au sens où la société l'entendait, mais un corps qui travaille, vieillit et prend soin d'autres corps. On pourrait y voir une déchéance ; le texte y voit l'apparition d'une norme nouvelle. C'est ce que Canguilhem nomme normativité : dans « Le normal et le pathologique », la santé n'est pas la conformité à une moyenne, c'est le pouvoir d'instituer de nouvelles normes, y compris dans des conditions dégradées. Ni le vivant ni le milieu ne sont normaux pris à part : c'est leur relation qui l'est. La narratrice est en ce sens en bonne santé : elle a fait d'un milieu impossible un milieu vivable.

Nemo offre la version sombre de cet effort. Il a voulu être pleinement ce qu'il est, « là seulement est l'indépendance ! là je ne reconnais pas de maîtres ! là je suis libre ! », et sa liberté s'achève dans la vengeance : l'éperonnage du navire de guerre, regardé jusqu'au bout, puis les sanglots devant le portrait d'une femme et de deux enfants. Devenir ce que l'on est peut signifier s'enfermer dans ce que l'on est. Aronnax trace une voie moyenne : ni Nemo ni Ned Land, il devient un savant transformé par ce qu'il a vu et rapporte à la surface un savoir et un récit. Son trajet est celui que Canguilhem assigne à la pensée : partie de la vie, séparée d'elle par l'analyse, elle y revient en la comprenant mieux.

\textbf{B. « À ce moment-là » : la mort est-elle la continuation logique ?}

Reste la phrase la plus contestable, qui commande la dernière : « Quand la mort arrivera, ne t'inquiète pas, c'est la continuation logique. » Giono naturalise la mort et lui confie un rôle : « à ce moment-là », chacun devient ce qu'il est, et plus la vie a été riche, plus ce qui retourne au tout est grand. Les trois œuvres se séparent ici nettement.

Verne semble donner raison à Giono, de la manière la plus littérale. À la fin de la première partie, un homme de l'équipage, blessé dans un choc dont Nemo ne parle qu'évasivement, est enseveli au fond de la mer, dans un cimetière de corail. Les polypes scelleront le corps ; les morts y dorment, dit Nemo, à l'abri des requins, et il ajoute : et des hommes. Le mort devient architecture vivante qui croît : image exacte de la continuation. Mais le dernier mot de Nemo trahit que ce tombeau est d'abord un refuge contre l'humanité, et le roman ne laisse pas à Nemo la sérénité qu'il prête à ses morts : il le montre en sanglots après le naufrage qu'il a provoqué, puis l'abandonne aux abords du Maelström, sans qu'Aronnax sache s'il a survécu.

Haushofer récuse la formule. À la fin de l'été, un inconnu surgit à l'alpage et tue à la hache le taureau, puis le chien Lynx ; la narratrice l'abat d'un coup de fusil. Aucun cycle ne se referme : il y a une perte sèche, une compagnie irremplaçable. Cette différence a une raison précise. Chez Giono, la mort est la mienne : le « ne t'inquiète pas » s'adresse à celui qui va mourir. Chez Haushofer, c'est la mort de l'autre, celle que le survivant doit porter, et aucune philosophie du cycle n'y répond.

Canguilhem permet de formuler ce refus. Si vivre, c'est poser des valeurs, la vie est polarisée : elle n'est pas indifférente entre vivre et mourir, elle préfère, et cette préférence est le fait biologique premier. Tenir la mort pour une « continuation logique », c'est adopter le point de vue de l'univers physique, où toutes les transformations s'équivalent, ce point de vue que Canguilhem refuse d'imposer à la biologie. Pour le vivant, la mort n'est pas une continuation : elle est la fin du centre de référence à partir duquel il y avait un monde. La monstruosité le rappelle : la vie peut se manquer, et une vie qui peut se manquer ne tient pas sa fin pour logique.

\textbf{C. La vraie richesse : ce qui reste quand tout est compté}

Que reste-t-il, une fois retirés le chiasme et la sérénité devant la mort ? L'essentiel, mais déplacé : la richesse n'est ni un stock ni une harmonie, c'est une relation entretenue, et ce que l'on « devient » tient moins au retour dans le tout qu'aux liens que l'on a noués.

\emph{Le Mur invisible} en fait la démonstration par les bêtes. Ce qui fait vivre la narratrice, ce n'est pas la beauté de l'alpage, c'est Bella qu'il faut traire, Lynx qui l'attend, la chatte, le jeune taureau. Elle ne peut pas se donner la mort tant que des bêtes dépendent d'elle. La richesse a ici un nom : la responsabilité. Le récit s'arrête faute de papier, et la narratrice annonce qu'elle ira nourrir la corneille blanche, que les autres corneilles rejettent à cause de sa couleur. Une bête inutile, qu'on nourrit parce qu'elle est là : voilà une richesse qui ne doit rien à un plan du monde et tout à une décision. Et ce qui demeure de la narratrice n'est pas un retour au tout, c'est un rapport écrit pour personne et pourtant lu.

Chez Verne, la même vérité affleure contre le héros. On retient de Nemo ses gestes de soin, le pêcheur arraché au requin, l'or rendu aux opprimés, le deuil de ses compagnons, et ce qui le perd est la rupture de tout lien avec ses semblables. La richesse de la mer n'a pas suffi : il fallait quelqu'un à qui la dire, et c'est pourquoi Aronnax est là. Le roman, récit ramené à la surface, est l'acte par lequel une richesse solitaire redevient partageable.

Canguilhem nomme la disposition qui rend cette richesse possible : une attention du connaissant à ce qu'il connaît. Sa réhabilitation du vitalisme demande au biologiste de ne pas oublier qu'il est lui-même un vivant, et de traiter son objet comme un centre de valeurs plutôt que comme un mécanisme muet. La connaissance de la vie devient ainsi une éthique : elle reconnaît à l'autre vivant sa manière propre d'habiter un monde. Le titre de Giono était juste, car la richesse se mesure à l'intensité des relations et non à ce que l'on détient ; il était trompeur, car le mot suggère encore un avoir, alors que ce dont il s'agit ne s'accumule pas et se perd dès qu'on cesse de l'entretenir.

\hypertarget{conclusion}{%
\subsection{Conclusion}\label{conclusion}}

Confronté aux trois œuvres, le propos de Giono sort profondément remanié. Son diagnostic tient : \emph{Vingt mille lieues sous les mers}, \emph{La Connaissance de la vie} et \emph{Le Mur invisible} s'écrivent depuis une séparation et retournent contre l'économie le mot de richesse, pour le rendre au vent, au lait, à la lumière et au corail. Sa métaphysique cède : rien n'a été fait pour nous, il n'existe pas de milieu absolu dont l'homme serait le destinataire, la finalité que sauve Canguilhem est intérieure au vivant et ne garantit rien, et l'on n'accède au monde qu'à travers une vitre, une faux ou un concept. Sa sérénité devant la mort, enfin, ne résiste pas à la mort des autres : le cimetière de corail de Verne apaise ce que la hache de l'inconnu, chez Haushofer, rend inconsolable.

Reste l'impératif, et c'est peut-être le vrai sujet. « Tâche seulement d'être le plus riche possible » : le verbe dit un travail, et ce travail est ce que les œuvres décrivent, la technique chez Verne, le concept chez Canguilhem, le soin chez Haushofer. Ce que l'on devient au terme n'est pas le retour au tout, mais la somme des liens que l'on a su instituer avec un monde qui ne nous attendait pas.

Une dernière remarque déplace la question. Le mur qui enferme la narratrice n'est pas venu de la nature ; il est sorti des mains des hommes. Chez Giono, des hommes nous privent des forêts ; chez Haushofer, ils privent les forêts de tout regard humain, sauf un. La question n'est peut-être plus de savoir ce dont on nous prive, mais ce que nous sommes encore capables de ne pas détruire.

\end{document}
